\exer{Modélisation géométrique}
\setcounter{exo}{0}
%Ecole de l'air PSI 2002

On considère le système suivant.
\begin{center}
\includegraphics[width=\linewidth]{973_01}%
\end{center}

Sur le dessin suivant, la position pour laquelle x est nul est représentée en trait fort et celle pour laquelle x est non nul en trait pointillé. On note lambda la longueur A prime B et lambda zéro la longueur AB.

\begin{center}
\includegraphics[width=\linewidth]{973_02}%
\end{center}

\question{Déterminer lambda et lambda zéro en fonction de R, L et x.}
\ifprof
\begin{corrige}
\end{corrige}
\else
\fi

\question{Déterminer la différence entre lambda zéro et lambda en fonction de R et de la différence angulaire entre theta c et theta.}
\ifprof
\begin{corrige}
\end{corrige}
\else
\fi

\question{En déduire une expression de x en fonction de la différence angulaire, de R et de L.}
\ifprof
\begin{corrige}
\end{corrige}
\else
\fi

\question{Les variations angulaires et le déplacement x sont faibles devant les autres grandeurs. Effectuer la linéarisation autour du point d’équilibre.}
\ifprof
\begin{corrige}
\end{corrige}
\else
\fi

\question{Déterminer la valeur de K reliant x à la variation angulaire.}
\ifprof
\begin{corrige}
\end{corrige}
\else
\fi

\begin{enumerate}
\item $\lambda_0 = \sqrt{L^2+R^2}$ et $\lambda = \sqrt{L^2 +\left(R-x^2\right)}$.
\item $2\left(\lambda_0 -\lambda\right)=R\Delta \theta$.
\item $x=R-\sqrt{R^2-R\sqrt{L^2+R^2}\Delta \theta+R^2\dfrac{\Delta \theta^2}{4}}$.
\item ...
\item $K=\dfrac{\sqrt{L^2+R^2}}{2}$.
\end{enumerate}

% BEGIN AUTO CORRIGE INCLUDE
\InclureCorrige{corrige.tex}
% END AUTO CORRIGE INCLUDE
